% TOP_PROBLEM: {"id": 30007269, "title": "Tame eigensheaves at several marked points", "category_id": 1, "category": {"id": 1, "name": "number_theory", "display_name": "Number Theory", "description": "Properties of integers, prime numbers, Diophantine equations.", "slug": "number-theory", "order_index": 1, "created_at": "2026-07-31T15:26:25.670Z"}, "status": "open"}
% FIRST_POSTED: 2026-10-06
% ATTEMPT_STATUS: solved
% OPENAI_CLAIM: {"schema_version":1,"source_repo":"https://github.com/openai/math","source_commit":"adc7f1241b42e322a6451854ab7e4b4c146bf78a","release_date":"2026-10-06","match":"full","resolution":"proved","independently_reviewed":false,"second_pass":"agreed","families":[{"family":"014","title":"Restricted geometric Langlands, global Arthur enhancements, and generic Ramanujan","subject":"Number theory","manuscripts":[{"title":"Tame Hecke Eigensheaves with Several Marked Points","pdf_path":"preprints/Tame-Hecke-Eigensheaves-with-Several-Marked-Points-October-5-2026/Tame-Hecke-Eigensheaves-with-Several-Marked-Points.pdf","readme_path":"preprints/Tame-Hecke-Eigensheaves-with-Several-Marked-Points-October-5-2026/README.md","date":"2026-10-05","theorem_ref":"intro:main, intro:eigen (sections/introduction.tex:40)","statement_sources":[{"path":"preprints/Tame-Hecke-Eigensheaves-with-Several-Marked-Points-October-5-2026/build/sections/introduction.tex","line":40,"label":"intro:main, intro:eigen"}]}],"lean":null,"reasoning_trace":null}]}
\documentclass[11pt]{article}
\usepackage[margin=25mm]{geometry}
\usepackage{fontspec}
\setmainfont{Latin Modern Roman}
\usepackage{amsmath,amssymb,mathtools}
\usepackage{unicode-math}
\setmathfont{Latin Modern Math}
\usepackage{enumitem}
\usepackage{xurl,hyperref,bookmark}
\hypersetup{colorlinks=true,urlcolor=blue}
\setcounter{secnumdepth}{0}
\setlength{\parindent}{0pt}
\setlength{\parskip}{5pt}
\setlength{\emergencystretch}{3em}
\begin{document}
\section*{Tame eigensheaves at several marked points}

\subsection{OpenAI's claim}
OpenAI presents the following claims relevant to this question, with manuscript numbers following paper order in Sources. In manuscript 1 (family 014; 2026-10-05), the asserted result is that for $\mathrm{SL}_n$ with $p>n$, dense tame parameters with arbitrary unipotent Jordan types at several marked points admit nonzero Borel-level perverse eigensheaves retaining tensor and fusion compatibility.

\subsection{Mathematical statement of the claimed result}
\paragraph{Manuscript 1.}
Let $n\ge2$, $k=\overline{\mathbb F}_q$ of characteristic $p>n$, $\ell\ne p$, and $E=\overline{\mathbb Q}_{\ell}$. Take a smooth projective geometrically connected $X_0/\mathbb F_q$ of genus at least two with distinct rational $x_1,\ldots,x_s$, $s\ge2$, and set $X=(X_0)_k$, $D=\sum_ix_i$, $U=X\setminus|D|$. A continuous Zariski-dense $\rho:\pi_1^{\rm et}(U)\to\operatorname{PGL}_n(L)$, $L/\mathbb Q_{\ell}$ finite, is assumed trivial on wild inertia and has
\[
 \rho(\gamma)=\exp(t_{\ell,i}(\gamma)N_i)\quad(\gamma\in I_{x_i}),
\]
with any nilpotent $N_i$, including zero. On $\mathcal A=\operatorname{Bun}_{\operatorname{SL}_n,B,D}(X)$ the claim supplies a nonzero locally constructible perverse geometric sheaf $M$ with singular support in the generically nilpotent Higgs cone with Borel-unipotent residues, and
\[
 \operatorname{Hecke}_{I,(V_i)}(M)\simeq M\boxtimes\mathop{\boxtimes}_{i\in I}(V_i)_{\rho}
\]
for all finite leg sets and dual-group representations. Eigenmaps are coherent for units, convolution, permutations and every collision diagonal, with relative Satake normalization $E[d_{\lambda}](d_{\lambda}/2)$ and no extra moving-leg shift.

\subsection{Outline of the argument}
For manuscript 1 (family 014; 2026-10-05), a curve with several separating nodes realizes all marked types simultaneously. Nilpotent cotangent dimension bounds and transverse modifications control specialization, after which commuting unipotent torus actions remove all flag enhancements without destroying the object.

\subsection{Scope relative to this problem}
The result supplies the conclusion in this problem with its displayed hypotheses and quantifiers. The broader statements and applications in the manuscript are kept separate from the particular question above. The supporting result is manuscript 1 (family 014; 2026-10-05).

\subsection{Formal verification}
No Lean formalization is listed at this snapshot for manuscript 1 (family 014; 2026-10-05).

This site has not executed the Lean code. Described statement coverage is not an independent proof check.

\subsection{Status caveat}
The manuscripts were produced by an internal OpenAI model. OpenAI cautions that results without formalization can contain mistakes. This site has not independently checked the proofs; the label records the scope of OpenAI's claim.

\subsection{Sources}
\begin{itemize}
\item OpenAI, \emph{Tame Hecke Eigensheaves with Several Marked Points}, 2026-10-05; citation key \texttt{OAI:Tame-Hecke-Eigensheaves-with-Several-Marked-Points-October-5-2026}. \href{https://github.com/openai/math/blob/adc7f1241b42e322a6451854ab7e4b4c146bf78a/preprints/Tame-Hecke-Eigensheaves-with-Several-Marked-Points-October-5-2026/Tame-Hecke-Eigensheaves-with-Several-Marked-Points.pdf}{PDF (GitHub)}; \href{https://github.com/openai/math/raw/adc7f1241b42e322a6451854ab7e4b4c146bf78a/preprints/Tame-Hecke-Eigensheaves-with-Several-Marked-Points-October-5-2026/Tame-Hecke-Eigensheaves-with-Several-Marked-Points.pdf}{PDF (direct)}; \href{https://github.com/openai/math/blob/adc7f1241b42e322a6451854ab7e4b4c146bf78a/preprints/Tame-Hecke-Eigensheaves-with-Several-Marked-Points-October-5-2026/README.md}{Manuscript page and citation}; \href{https://github.com/openai/math/blob/main/preprints/Tame-Hecke-Eigensheaves-with-Several-Marked-Points-October-5-2026/Tame-Hecke-Eigensheaves-with-Several-Marked-Points.pdf}{Latest version}; Statement source: \href{https://github.com/openai/math/blob/adc7f1241b42e322a6451854ab7e4b4c146bf78a/preprints/Tame-Hecke-Eigensheaves-with-Several-Marked-Points-October-5-2026/build/sections/introduction.tex\#L40}{intro:main, intro:eigen}.
\item \href{https://github.com/openai/math/blob/adc7f1241b42e322a6451854ab7e4b4c146bf78a/CONTENTS.md}{OpenAI catalogue entry}.
\end{itemize}
\end{document}
